% GENERATED FILE. Edit canonical lessons, then run npm run generate:books.
% canonical-source-hash: fnv1a64:26fe321a1f2195ec
% canonical-lessons: FR-C09-juillet, FR-C09-aout, FR-C09-septembre, FR-C09-octobre, FR-C09-novembre, FR-C09-decembre, FR-C09-mois-practice-2

\chapter{The Months: Caesars and Counting}
\label{ch:months-caesars}
\hlchaptermodality{10}

\begin{chapteropening}
\textbf{By the end of this chapter:} \emph{I can name the last six months in French and explain why four of them are counting wrong.}

Two men put themselves in a calendar of gods, and the four months after them have been miscounting ever since. juillet is Julius Caesar and aout is the emperor Augustus; wedging those two into the middle is what pushed septembre to decembre two places out of place.
\end{chapteropening}

\section[juillet]{juillet — July}

\label{lesson:FR-C09-juillet}



\begin{quote}
Six months named for gods and a festival. The next two are named for \textbf{men} — and both men put their own months into the calendar themselves.
\end{quote}

\begin{sounds}
\begin{itemize}
\raggedright
  \item \texttt{ill-mouille} — the \textbf{ill} in the middle is not an \emph{l} sound at all. It is a \emph{y}, as in English \emph{yes}. This spelling is everywhere in French and it catches every beginner once.
  \item With the soft \textbf{j} of \textbf{juin}: \textbf{juillet} = \emph{zhwee-YAY}.
\end{itemize}
\end{sounds}

\begin{cousinweb}
\textbf{juillet} is Latin \textbf{Julius}, and the Julius in question is \textbf{Julius Caesar}.

He did not merely lend his name to an existing month. He \textbf{reformed the entire calendar} — the Julian calendar is his, and it ran Europe for sixteen centuries — and while rearranging it he renamed one month after himself.

English \textbf{July} is the same name, worn down a different way from the same \emph{Julius}.
\end{cousinweb}

\subsection*{Your turn}
Say these aloud:

\begin{itemize}
\raggedright
  \item \textquotedblleft{}juillet\textquotedblright{} — \emph{zhwee-YAY}; the \emph{ill} is a \textbf{y}, never an \emph{l}
  \item \textquotedblleft{}juin, juillet\textquotedblright{} — both open with the soft \emph{j}, and they are easy to confuse; say them back to back until they separate
  \item the man — \textquotedblleft{}Julius Caesar, who rebuilt the calendar and kept a month\textquotedblright{}
\end{itemize}

\begin{tcolorbox}[breakable,colback=teal!4,colframe=teal!35!black,title={Before you move on}]
How do you say July? (\textbf{juillet}, \emph{zhwee-YAY}.) What sound does the \textbf{ill} make? (A \textbf{y}, as in \emph{yes} — never an \emph{l}.) Who is it named for, and what did he do to the calendar? (\textbf{Julius Caesar} — he \textbf{reformed} it, and named this month for himself.) Next: the emperor who copied him.
\end{tcolorbox}
\section[août]{août — August}

\label{lesson:FR-C09-aout}



\begin{quote}
The most violently worn-down word in the French calendar. Four letters on the page, and barely two sounds in the mouth.
\end{quote}

\begin{sounds}
\begin{itemize}
\raggedright
  \item \textbf{août} = \emph{oot}, and in much of France simply \emph{oo}. The \textbf{a} is not sounded, and the final \textbf{t} is optional.
  \item \texttt{circumflex} — the little hat on the \textbf{û} is a \textbf{gravestone}. It marks a letter that used to be there and is gone: older French wrote \emph{aoust}, and when the \textbf{s} fell silent the hat took its place.
\end{itemize}

You have met that hat before, on \textbf{hôpital} and \textbf{forêt} in English's \emph{hospital} and \emph{forest}. Wherever you see a circumflex, try putting an \textbf{s} after the vowel and see whether an English word appears.
\end{sounds}

\begin{cousinweb}
\textbf{août} is Latin \textbf{Augustus}, the first Roman emperor.

Julius Caesar took a month. His heir Augustus, not to be outdone, \textbf{took the next one}. Two men, two months, inserted into the middle of the year by the two people powerful enough to do it.

That insertion is the thing you already know the consequence of. Wedging two months into the middle is what pushed the counting months out of place — which is why the ninth month still says \textquotedblleft{}seven.\textquotedblright{}
\end{cousinweb}

\begin{culture}
Ten of the twelve months are a god, a festival or a number. \textbf{Two are people.}

Julius and Augustus put themselves among the gods, in the most public place available — the calendar everybody has to say out loud. Two thousand years later, France still says both their names every summer, and nobody thinks of either man while doing it. That is what winning looks like at scale.
\end{culture}

\subsection*{Your turn}
Say these aloud:

\begin{itemize}
\raggedright
  \item \textquotedblleft{}août\textquotedblright{} — \emph{oot}, and hear how little of the spelling survives
  \item \textquotedblleft{}juillet, août\textquotedblright{} — the two men, back to back
  \item the hat — \textquotedblleft{}aoust lost its \emph{s}, and the circumflex marks the grave\textquotedblright{}
\end{itemize}

\begin{tcolorbox}[breakable,colback=teal!4,colframe=teal!35!black,title={Before you move on}]
How do you say August? (\textbf{août} — \emph{oot}, or simply \emph{oo}.) What does the circumflex mark? (A \textbf{lost letter} — here an \emph{s}, from older \emph{aoust}.) Who are the two months named for people, and what did those people do? (\textbf{Julius Caesar} and the emperor \textbf{Augustus} — each inserted a month and kept it.) What did that insertion do to the rest of the calendar? (\textbf{Pushed the counting months two places out.}) Next: the four that are still counting.
\end{tcolorbox}
\section[septembre]{septembre — September}

\label{lesson:FR-C09-septembre}



\begin{quote}
You met this word while counting, as the proof that \emph{sept} means seven. Now take it as a month in its own right.
\end{quote}

\begin{sounds}
\begin{itemize}
\raggedright
  \item \textbf{septembre} = \emph{sep-TAHM-bruh}. The \textbf{em} in the middle is \textbf{nasal}, and the closing \textbf{-re} is barely a syllable.
  \item All four of the counting months end this same way, so learning it once buys you four.
\end{itemize}
\end{sounds}

\begin{cousinweb}
\textbf{septembre} is Latin \textbf{september}, \textquotedblleft{}the \textbf{seventh}\textquotedblright{} — from \textbf{septem}, the seven you can already count with.

You know why it sits ninth: the old Roman year began in March, and then Julius and Augustus wedged their two months in. The name never caught up. What is new here is the \textbf{word}: \emph{septembre} is now yours as a month, not only as evidence about a number.

English \textbf{September} is the identical Latin word, carrying the identical mistake.
\end{cousinweb}

\subsection*{Your turn}
Say these aloud:

\begin{itemize}
\raggedright
  \item \textquotedblleft{}septembre\textquotedblright{} — \emph{sep-TAHM-bruh}, the \emph{em} through the nose
  \item \textquotedblleft{}sept, septembre\textquotedblright{} — the number and the month it hides in
  \item \textquotedblleft{}août, septembre\textquotedblright{} — the emperor, then the miscount he caused
\end{itemize}

\begin{tcolorbox}[breakable,colback=teal!4,colframe=teal!35!black,title={Before you move on}]
How do you say September? (\textbf{septembre}, \emph{sep-TAHM-bruh}.) What number is inside it? (\textbf{sept}, seven — Latin \emph{septem}.) Which month is it actually? (The \textbf{ninth}.) Next: the eighth month, which is the tenth.
\end{tcolorbox}
\section[octobre]{octobre — October}

\label{lesson:FR-C09-octobre}



\begin{quote}
Of the four counting months, this is the one that rescues a number you could no longer see.
\end{quote}

\begin{sounds}
\begin{itemize}
\raggedright
  \item \textbf{octobre} = \emph{ok-TOH-bruh}. The \textbf{o} sounds are both closer to English \emph{taught} than to \emph{toe}.
  \item Same \emph{-bruh} ending as \textbf{septembre}.
\end{itemize}
\end{sounds}

\begin{cousinweb}
\textbf{octobre} is Latin \textbf{october}, \textquotedblleft{}the \textbf{eighth},\textquotedblright{} from \textbf{octō}.

Here is why this word matters more than the other three. When you learned to count, \emph{octō} had been battered past recognition on its way to \textbf{huit} — the \emph{c} gone, the vowels softened, a written \emph{h} propped on the front. The number kept none of its ancestry.

\textbf{The month kept all of it.} \emph{Oct-} stands there in plain sight, exactly as Latin left it, and English \textbf{octave}, \textbf{octagon} and \textbf{octopus} stand beside it. The eight vanished from the number and survived in the calendar.
\end{cousinweb}

\subsection*{Your turn}
Say these aloud:

\begin{itemize}
\raggedright
  \item \textquotedblleft{}octobre\textquotedblright{} — \emph{ok-TOH-bruh}
  \item the contrast — \textquotedblleft{}huit, octobre\textquotedblright{}: the same eight, unrecognisable and intact
  \item \textquotedblleft{}septembre, octobre\textquotedblright{} — the pair
\end{itemize}

\begin{tcolorbox}[breakable,colback=teal!4,colframe=teal!35!black,title={Before you move on}]
How do you say October? (\textbf{octobre}, \emph{ok-TOH-bruh}.) What number is inside it? (\textbf{octō}, eight.) Why is the eight easier to see here than in the number \textbf{huit}? (\emph{Huit} wore its Latin away; the \textbf{month kept it}.) Next: the ninth, which is the eleventh.
\end{tcolorbox}
\section[novembre]{novembre — November}

\label{lesson:FR-C09-novembre}



\begin{quote}
The third of the four, and by now you can predict both the shape and the arithmetic.
\end{quote}

\begin{sounds}
\begin{itemize}
\raggedright
  \item \textbf{novembre} = \emph{no-VAHM-bruh}, the \textbf{em} nasal and the ending the same \emph{-bruh} as the last two.
  \item Three months in a row have now ended identically. That is the pattern paying out.
\end{itemize}
\end{sounds}

\begin{cousinweb}
\textbf{novembre} is Latin \textbf{november}, \textquotedblleft{}the \textbf{ninth},\textquotedblright{} from \textbf{novem} — the nine you count with as \textbf{neuf}.

Unlike \emph{huit}, \textbf{neuf} still shows its Latin: \emph{novem} to \emph{neuf} is a short walk, and the \emph{nov-} is right there in both. English keeps the same root in \textbf{novena}, a devotion of nine days.

Be careful of one false friend: Latin \textbf{novus} means \textquotedblleft{}new,\textquotedblright{} and English \textbf{novel} and \textbf{novelty} come from \emph{that}, not from \emph{novem}. Two similar roots, two different ideas. \textbf{November is a count, not a newness.}
\end{cousinweb}

\subsection*{Your turn}
Say these aloud:

\begin{itemize}
\raggedright
  \item \textquotedblleft{}novembre\textquotedblright{} — \emph{no-VAHM-bruh}
  \item \textquotedblleft{}neuf, novembre\textquotedblright{} — the nine in both
  \item \textquotedblleft{}septembre, octobre, novembre\textquotedblright{} — three in a row, one ending
\end{itemize}

\begin{tcolorbox}[breakable,colback=teal!4,colframe=teal!35!black,title={Before you move on}]
How do you say November? (\textbf{novembre}, \emph{no-VAHM-bruh}.) What number is inside it? (\textbf{novem}, nine — the \emph{neuf} you count with.) Does it have anything to do with \textquotedblleft{}new\textquotedblright{}? (\textbf{No} — that is Latin \emph{novus}, a different root.) Next: the last of them, and the biggest number.
\end{tcolorbox}
\section[décembre]{décembre — December}

\label{lesson:FR-C09-decembre}



\begin{quote}
The twelfth month, and it says \textbf{ten}. This is the last and largest of the four miscounts, and after it the year closes.
\end{quote}

\begin{sounds}
\begin{itemize}
\raggedright
  \item \textbf{décembre} = \emph{day-SAHM-bruh}. The \textbf{é} is the bright \emph{ay} of \textbf{février}, the \textbf{c} before \emph{e} is an \emph{s}, and the \textbf{em} is nasal.
  \item The familiar \emph{-bruh} closes it for the fourth and final time.
\end{itemize}
\end{sounds}

\begin{cousinweb}
\textbf{décembre} is Latin \textbf{december}, \textquotedblleft{}the \textbf{tenth},\textquotedblright{} from \textbf{decem}.

That root is one of the most productive you have met. It gave you the number \textbf{dix}; it gives English \textbf{decimal}, \textbf{decade} and \textbf{decimate}; and it reached an American pocket as the \textbf{dime}, through Old French \emph{disme} from Latin \emph{decima}, \textquotedblleft{}a tenth part.\textquotedblright{}

So every December, in French and English alike, you say the Roman word for ten and mean the twelfth month. Spanish carries the identical fault: \emph{diciembre} is \emph{decem} too.
\end{cousinweb}

\subsection*{Your turn}
Say these aloud:

\begin{itemize}
\raggedright
  \item \textquotedblleft{}décembre\textquotedblright{} — \emph{day-SAHM-bruh}
  \item \textquotedblleft{}dix, décembre, dime\textquotedblright{} — one Latin \emph{decem} in three places
  \item all four — \textquotedblleft{}septembre, octobre, novembre, décembre\textquotedblright{}
\end{itemize}

\begin{tcolorbox}[breakable,colback=teal!4,colframe=teal!35!black,title={Before you move on}]
How do you say December? (\textbf{décembre}, \emph{day-SAHM-bruh}.) What number is inside it? (\textbf{decem}, ten — the \emph{dix} you count with.) Name two English words from the same root. (\emph{Decimal}, \emph{decade}, \emph{decimate} — or the \textbf{dime}.) Next: all twelve, together.
\end{tcolorbox}
\section[Practice]{Practice — the whole year}

\label{lesson:FR-C09-mois-practice-2}



\begin{quote}
Twelve months, six at a time. Run the whole year now, and hear where it breaks.
\end{quote}

\subsection*{Your turn: the twelve, in order}
Read down the left, then cover the right and read back:

\noindent
\begin{tabularx}{\linewidth}{@{}>{\raggedright\arraybackslash}X>{\raggedright\arraybackslash}X@{}}
\toprule
\textbf{French} & \textbf{English} \\
\midrule
\textbf{janvier} & January \\
\textbf{février} & February \\
\textbf{mars} & March \\
\textbf{avril} & April \\
\textbf{mai} & May \\
\textbf{juin} & June \\
\textbf{juillet} & July \\
\textbf{août} & August \\
\textbf{septembre} & September \\
\textbf{octobre} & October \\
\textbf{novembre} & November \\
\textbf{décembre} & December \\
\bottomrule
\end{tabularx}

\emph{Twice through:} the twelve months, in order

Then start from \textbf{juillet} and go round to \emph{juin}.

\subsection*{Your turn: the two hard pairs}
Two pairs cause almost all the trouble. Drill them alone:

\begin{itemize}
\raggedright
  \item \textbf{juin} / \textbf{juillet} — \emph{zhwa(n)} against \emph{zhwee-YAY}. Both open with the soft \emph{j} and neither is long.
  \item \textbf{mars} / \textbf{mai} — \emph{marss} against \emph{may}. The two short ones, and the two most easily swallowed.
\end{itemize}

\begin{grammarlens}[title={What you've built this chapter}]
\begin{itemize}
\raggedright
  \item \textbf{twelve month names}, and the year runs end to end.
  \item \textbf{two men in a calendar of gods} — \emph{juillet} for Julius Caesar, \emph{août} for the emperor Augustus. Each inserted a month; each kept it.
  \item \textbf{why the last four miscount} — wedging those two into the middle pushed \emph{septembre} to \emph{décembre} two places along, and their Latin names never caught up.
  \item \textbf{a number rescued} — \emph{octō} is unrecognisable inside \textbf{huit} and perfectly legible inside \textbf{octobre}. When a word wears down, look for a longer relative that did not.
\end{itemize}
\end{grammarlens}

\begin{tcolorbox}[breakable,colback=teal!4,colframe=teal!35!black,title={Before you move on}]
Give the twelve months without stopping. Now start from \textbf{septembre}. Which two are named for people, and who? (\textbf{juillet}, Julius Caesar; \textbf{août}, the emperor Augustus.) What did those two insertions do? (Pushed the \textbf{counting months two places out}.) Which month shows a number that its own French numeral has lost? (\textbf{octobre} — the \emph{octō} that vanished from \emph{huit}.)
\end{tcolorbox}
